%%%==========================================================================%%%
%%%              Packages and such like added by Norman Dunbar               %%%
%%%==========================================================================%%%

\usepackage{nameref}% To use \nameref{...}
\usepackage{wrapfig}% To wrap text around figures.
\usepackage{listings}% Guess!
\usepackage[toc,page]{appendix}% Gets a ``blank'' Appendix page.
\usepackage{longtable}% For tables extending over a page.
\usepackage[obeyspaces]{url}% Allow spaces in \path{}.
\usepackage{xurl}% Ensure \URLs break at /.:*-~'" characters.
\usepackage{xcolor}% Required for the setup of the listings.
\usepackage[LGR,T1]{fontenc}% To get Ohms symbols.
\usepackage{array}%To get table's centred and wrapped.
\usepackage{courier}% For listings, the default is dire!
\usepackage[unicode=true,
    bookmarks=true,
    bookmarksnumbered=false,
    bookmarksopen=false,
    breaklinks=true,
    pdfborder={0 0 0},
    pdfborderstyle={},
    backref=false,
    colorlinks=true]{hyperref}
\usepackage{multicol}
\usepackage{multirow}

\urlstyle{same}% Ensure URLs are in the same font as the text. Not in monotype!

% Macros.
\newcommand{\ohms}{\textgreek{W}}%Ohms symbol.
%\newcommand{\app}[1]{\emph{#1}}%Application names.
\newcommand{\bin}[1]{0b{#1}}%Binary numbers.
\newcommand{\hexa}[1]{0x\uppercase{{#1}}}%Hexadecimal numbers.
\newcommand{\varb}[1]{\texttt{#1}}%Variable names.
\newcommand{\macro}[1]{\texttt{#1}}%Defined macro names.
\newcommand{\code}[1]{\texttt{#1}}%Inline code segments.
\newcommand{\func}[1]{\texttt{#1}}%Procedure and function names.
\newcommand{\pin}[1]{\texttt{#1}}%Arduino/AVR pin names.
\newcommand{\externC}{\code{extern \textquotedbl{}C\textquotedbl}}
\newcommand{\shl}{\texttt{<}\texttt{<}~}% << in inline code.
\newcommand{\shr}{\texttt{>}\texttt{>}~}% >> in inline code.
\newcommand{\atmega}{ATmega328P}% Get capitalisation correct!
\newcommand{\gpio}[1]{\texttt{GPIO~#1}}% GPIO pin names.
\newcommand{\reg}[1]{\texttt{\uppercase{{#1}}}}%Hexadecimal numbers.
\newcommand{\lyxarrow}{$\rightarrow$}
\newcommand{\chapref}[1]{Chapter \ref{#1}, \nameref{#1}}
\newcommand{\sectref}[1]{Section \ref{#1}, \nameref{#1}}
\newcommand{\class}[1]{\texttt{#1}}% C++ class names.

%Macros specific to this book.
\newcommand{\sdkversion}{2.3.0}%Current version of the SDK.


%%% Setup Hyperref configuration %%%
\hypersetup{pdftitle={A Beginner's Guide to Github},
    pdfauthor={Norman Dunbar},
    pdfsubject={A Guide on how to use GitHub.},
    pdfkeywords={github}
}


%%% Lyx specific additions %%%
\makeatletter

\DeclareRobustCommand{\greektext}{%
    \fontencoding{LGR}\selectfont\def\encodingdefault{LGR}}
\DeclareRobustCommand{\textgreek}[1]{\leavevmode{\greektext #1}}
\ProvideTextCommand{\~}{LGR}[1]{\char126#1}

\makeatother



\definecolor{backcolour}{rgb}{1.0, 0.95, 0.96}
\definecolor{codeGreen}{rgb}{0,0.6,0}
\definecolor{codeGray}{rgb}{0.5,0.5,0.5}
\definecolor{codePurple}{rgb}{0.58,0,0.82}
\definecolor{codeBlue}{rgb}{0,0,0.9}


%%% Settings for listings %%%
%\lstset{
\lstdefinestyle{myStyle}{
    language=C++,
    captionpos=b,
    %backgroundcolor=\color{backcolour},
    commentstyle=\itshape\color{codeGreen!},
    stringstyle=\color{codePurple},
    keywordstyle={\textbf{\color{black}}},
    identifierstyle=\color{codeBlue},
    breaklines=true,
    numbers=left,
    numberstyle=\color{codeGray},
    emptylines=*5,
    tabsize=4,
    basicstyle={\footnotesize\ttfamily}, %%% Overridden by SVMono class. :(
    %basewidth={0.55em}, %%% Needed to fit listings on page %%%
    showspaces=false,
    showstringspaces=false,
    keepspaces=true,
    frame=single,
    escapeinside={@*}{*@}, % Use to add LaTeX to a code lstlisting,
                           % \textbf{xxx\_yyy} --- you need to escape stuff though.
}

%%%==========================================================================%%%
%%%                End of Norman Dunbar's additions                          %%%
%%%==========================================================================%%%
